% Nature Cardiovascular Research Letter: original neutral layout sample.
% Source design measurements only. No source paper text, art or identifying metadata.
\PassOptionsToPackage{no-math}{fontspec}
\documentclass[UTF8,fontset=none,scheme=plain]{ctexart}
\usepackage[paperwidth=595.276bp,paperheight=790.866bp,left=39.685bp,right=34.016bp,
  top=49bp,bottom=40bp,headheight=15bp,headsep=10bp,footskip=15bp]{geometry}
% Shared dependencies and neutral material. Compile with XeLaTeX.
% No publisher logos, paper text, data, figures, or download metadata are included.
\usepackage{amsmath,amssymb,graphicx,xcolor,tikz,fancyhdr,needspace,xurl,hyperref,balance,etoolbox}
% The legacy balance package omits the first split box's final baseline depth.
% Reserve that measured depth rather than hiding a CJK overfull-box warning.
\makeatletter
\patchcmd{\@BAlancecol}{\setbox\@leftcolumn=\vbox to \@tempdima}{\advance\@tempdima by \dp1\relax\setbox\@leftcolumn=\vbox to \@tempdima}{}{\PackageWarning{cns-template}{Balance depth patch unavailable; inspect final columns}}
\makeatother
\usetikzlibrary{arrows.meta,positioning}
\usepackage[protrusion=true,expansion=false]{microtype}
\IfFontExistsTF{Noto Serif CJK SC}{%
  \setCJKmainfont{Noto Serif CJK SC}%
}{%
  \setCJKmainfont[BoldFont=FandolSong-Bold.otf,ItalicFont=FandolKai-Regular.otf]{FandolSong-Regular.otf}%
}
\IfFontExistsTF{Noto Sans CJK SC}{%
  \setCJKsansfont{Noto Sans CJK SC}%
}{%
  \setCJKsansfont[BoldFont=FandolHei-Bold.otf]{FandolHei-Regular.otf}%
}
\setmainfont{TeX Gyre Termes}\setsansfont{TeX Gyre Heros}
\setmonofont{TeX Gyre Cursor}
\definecolor{JournalBlue}{HTML}{0097C9}
\definecolor{JournalRed}{HTML}{B83443}
\definecolor{ScienceRed}{HTML}{CD1021}
\definecolor{CiteBlue}{HTML}{0000EE}
\hypersetup{colorlinks=true,linkcolor=JournalBlue,urlcolor=JournalBlue,
  pdftitle={CNS论文汉化：中性版式示例},pdfauthor={示例模板；非期刊官方模板},
  pdfsubject={Neutral source-inspired Chinese layout demonstration; no real research claims}}
\urlstyle{same}\raggedbottom
\setlength{\parindent}{1em}\setlength{\parskip}{0pt}
\setlength{\emergencystretch}{2em}\tolerance=2200
\XeTeXlinebreaklocale "zh"\XeTeXlinebreakskip=0pt plus .1pt
\pagestyle{fancy}\fancyhf{}
\renewcommand{\headrulewidth}{0pt}\renewcommand{\footrulewidth}{0pt}
\newcommand{\PaperTitle}{中文题名占位：从组织结构到证据解释}
\newcommand{\SampleAuthors}{Author A\textsuperscript{1}, Author B\textsuperscript{2,*}}
\newcommand{\SampleAffiliations}{\textsuperscript{1}作者单位一（示例占位）。\quad
  \textsuperscript{2}作者单位二（示例占位）。\par
  *通信作者信息待填。DOI、收稿和出版日期均应从目标原文核实。}
\newcommand{\SampleAbstract}{本段是中性摘要占位文字，不陈述任何真实研究发现。它演示中文题名、作者信息、摘要和双栏正文之间的层级关系。实际使用时，应完整替换为目标论文的忠实中文译文，并保留研究对象、方法、结果和结论的限定条件。图注、参考文献、作者声明和来源信息均属于独立的覆盖项。模板中的图形仅说明排版结构，不表示生物学机制、实验结果或统计关系。}
\newcommand{\DemoCaption}{\textbf{图1. 中性证据链示意（占位图，无实验数据）。}
  \textbf{A}，输入节点表示需要核对的原文材料。
  \textbf{B}，处理节点表示逐段转换及复核。
  \textbf{C}，输出节点表示正文、图注和书目等独立交付部分。
  箭头仅表示本示例的阅读顺序，不代表因果关系、效应方向或数量变化。
  该图由模板自行绘制，所有节点均为文字占位；无样本量、误差线、统计检验或比例尺。
  实际论文图注应逐项保留原稿报告的样本量、单位、检验方法、限定及图形署名。}
\newcommand{\DemoVector}{%
 \resizebox{\linewidth}{!}{%
 \begin{tikzpicture}[font=\sffamily\small,node distance=18mm,
   box/.style={draw=black!70,rounded corners=2pt,fill=black!3,minimum width=26mm,minimum height=16mm,align=center},
   flow/.style={-{Stealth[length=2.5mm]},line width=.7pt}]
  \node[box] (a) {\textbf{A}\quad 输入\\原文材料};
  \node[box,right=of a] (b) {\textbf{B}\quad 处理\\转换与复核};
  \node[box,right=of b] (c) {\textbf{C}\quad 输出\\中文译稿};
  \draw[flow] (a)--(b);\draw[flow] (b)--(c);
  \node[below=11mm of a,align=center] {标题与摘要\\正文与图表};
  \node[below=11mm of b,align=center] {保留限定\\检查单位};
  \node[below=11mm of c,align=center] {可追溯引文\\明确原文信息};
  \node[below=29mm of b,font=\sffamily\footnotesize,text=black!65] {示意占位；不包含真实实验数据};
 \end{tikzpicture}}%
}
\newcommand{\SampleBibliography}{%
 \section*{参考文献（示例占位）}
 {\fontsize{7.6bp}{10bp}\selectfont\setlength{\parindent}{0pt}\interlinepenalty=10000
 \hypertarget{ref:1}{}1. Author A. (20XX). Example reference title. \textit{Journal placeholder}. 卷、页码与DOI待填。此条仅供演示引文跳转。\par\smallskip
 \hypertarget{ref:2}{}2. Author B. (20YY). Another example reference title. \textit{Journal placeholder}. 出版信息待填；不要把占位条目带入正式译稿。\par}
}
\newcommand{\SampleSourceInfo}{%
 \section*{原文信息}
 以下字段全部为示例占位：作者；英文原标题；期刊；年份；卷期；页码或文章号；DOI或来源链接。实际使用时必须逐项替换并核对。\par
 本文件是\JournalName 风格的非官方中文译稿样板，依据一次已核验文章版式进行结构化改编，不代表该期刊所有年代或文体的统一标准，也不保证像素级复刻。无真实论文内容或期刊标识图像。\par
}
% Deliberately neutral demo prose, not translated from a published paper.
\newcommand{\SampleOpening}{%
 \section*{引言：示例正文与引文}
 这是一段用于检查双栏阅读节奏的中性示例文字。段落中的术语、缩写和标点需要在实际项目中与源文保持一致。中文译稿应呈现原论文的完整论述，而不是把多段证据压缩成一段摘要。这里的引文只链接到本文件末尾明确标为占位的参考文献\SampleCiteOne。\par
 若原句仅表示相关、可能或提示，译文也应保留相同的证据强度。数字、单位、正负号、上下标和范围端点应逐项核对。原文存在内部矛盾时，应分别保留相应位置的陈述，再用简短译注说明；不能根据常识把不同的数字静默统一。\par
 \subsection*{内容层级与跨页衔接}
 此小标题演示二级层级。实际论文可能采用不同的小标题密度，模板不会替代对原文结构的识别。双栏文章需按左栏、右栏和跨页续句的顺序读取。若段落在页末中断，应在完成整句核对后接入下一页，避免把图中文字误当正文，或重复翻译同一段。\par
 图表具有独立的信息范围。正文引用到某幅图，并不意味着其图注可以省略。图注中的每个分面、比较对象、时间点、样本量及统计说明都应对应源文。英文图内标签可以保留，并在图旁或文末提供中文释义；实际采用哪种安排，需要依据图形尺寸和译文长度判断\SampleCiteTwo。\par
 \subsection*{来源与可维护数据}
 本段继续演示正文密度。可编辑文本和图注数据应与最终版面分开维护，以便修订时准确定位变化。若只调整版式，应能证明正文、图注和书目内容没有被意外改写。更换字体或纸张后，须重新检查断行、图注与图片的相对位置以及页脚可读性。\par
 文内引用与书目之间应有明确对应关系。这里的两条演示书目没有真实DOI，也没有虚构的实验信息。使用模板时可增删条目，但必须同时更新编号、链接目标和引用位置。页码只表示中文译稿自己的分页，不用于冒充原出版页码。\par
 \subsection*{结果叙述的演示占位}
 本节的文字仅用于呈现长段落在双栏页面中的流动，不报告真实结果。实际翻译应保留观察对象、实验条件和结局指标之间的关系。例如，原文分别描述体外实验与动物实验时，不能把两类结果合并为同一层面的结论；研究使用的对照、时间顺序和分析范围也必须保留。图表中的比较范围不应扩大为对整个人群或全部疾病阶段的判断。\par
 论述中的转折与限制同样属于正文。若研究先提出一种解释，再讨论另一种可能，译文应沿用这一推进顺序。作者明确写出的不确定性应能在中文中被识别；如果尚未确定某个现象的方向或机制，就不能改成已经建立的因果关系。对于术语近似但研究层级不同的概念，可在首次出现时保留缩写，后文保持统一。\par
 \subsection*{方法信息与可复核性}
 方法部分往往包含密集数字和专有名词。翻译时须区分浓度、总量、频率和处理时长，不能因为中文句式调整而改变参数的归属。试剂名称、软件版本和数据访问标识通常保留原样；上下限、排除规则和显著性标准则需要逐字符核对。这里没有提供任何真实实验参数，避免把模板占位误认为可以直接执行的方法。\par
 表格与补充材料的引用也需要保留原层级。正文对补充图的指向并不意味着模板中已经附有该补充图；若源文件只包含主文，应如实标注所交付的范围。图表编号、附录顺序和参考文献数量都应由目标文章决定，本示例中的结构不构成固定配额，也不要求所有论文采用同样的章节。\par
}
\newcommand{\SampleClosing}{%
 \section*{讨论：版式适配与范围}
 上方矢量图是模板原创的排版占位图，不来自任何论文。图形与完整演示图注被放在同一页，便于读者直接核对分面说明。真实项目应使用经过授权或用户提供的原图素材，并保留图注署名与许可信息。尺寸不足时，可适度缩放或调整栏宽，不以删除图注解决空间问题。\par
 模板仅记录一次已核验文章的可观察设计特征。即使文章出自同一期刊，不同年代、研究论文与综述、校正清样与正式版本也可能采用不同的页眉和引文形式。应以用户提供的实际PDF为准，不能仅凭期刊名称推断配色、字号或引用方式\SampleCiteOne。\par
 \subsection*{演示检查项}
 本样板用于展示可编译性、中文字形、标题层级、跨栏图注、页眉页脚和内部链接。图中没有实际数据，因此不应被当作科学插图引用。替换内容后，还需检查所有页的渲染结果，确认无截字、碰撞、异常空白和遗漏段落。\par
 \SampleBibliography
 \SampleSourceInfo
}


\usepackage{multicol,booktabs,tabularx}
\definecolor{LetterPink}{rgb}{.9299916,.0649882,.3929961}
\definecolor{LetterBlue}{HTML}{0069A5}
\definecolor{ReportRed}{HTML}{ED1C24}
\hypersetup{linkcolor=LetterBlue,urlcolor=LetterBlue,
  pdftitle={Nature Cardiovascular Research：研究简报中性样板}}
\setlength{\columnsep}{12bp}
\setlength{\multicolsep}{5bp plus 2bp minus 2bp}
\setlength{\parindent}{2em}
\renewcommand{\normalsize}{\fontsize{9.15bp}{12.4bp}\selectfont}
\ctexset{section={format=\sffamily\bfseries\fontsize{12bp}{15bp}\selectfont,beforeskip=10bp,afterskip=4bp},
  subsection={format=\sffamily\bfseries\fontsize{9.6bp}{12.4bp}\selectfont,beforeskip=7bp,afterskip=2bp}}
\newcommand{\LetterType}{\sffamily\bfseries\fontsize{9bp}{11bp}\selectfont 研究简报\enspace Letter}
\newcommand{\LetterIdentifier}{\sffamily\fontsize{8bp}{10bp}\selectfont\textcolor{LetterBlue}{DOI：按所给原稿填写}}
\newcommand{\LetterCite}[1]{\textsuperscript{\hyperlink{ref:#1}{#1}}}
\newcommand{\ReportHeading}[1]{\par\medskip\noindent{\sffamily\bfseries\textcolor{ReportRed}{#1}}\par
  \vspace{2bp}{\color{ReportRed}\hrule height .5bp}\vspace{5bp}}
\fancyhead[L]{\LetterType}
\fancyhead[R]{\LetterIdentifier}
\fancyfoot[L]{\sffamily\fontsize{7.2bp}{8.6bp}\selectfont
  \textcolor{LetterBlue}{Nature Cardiovascular Research}\enspace｜\enspace 卷期及原刊页码待填}
\fancyfoot[R]{\sffamily\fontsize{7.2bp}{8.6bp}\selectfont 非官方中文样板\enspace｜\enspace 中文译页\ \thepage}
\renewcommand{\headrulewidth}{.35bp}
\renewcommand{\footrulewidth}{.35bp}
\fancypagestyle{letteropening}{\fancyhead{}\renewcommand{\headrulewidth}{0pt}}
\fancypagestyle{letterreport}{\fancyhead[L]{\sffamily\bfseries\fontsize{9bp}{11bp}\selectfont 报告摘要（中性演示）}}

\begin{document}\normalsize
\thispagestyle{letteropening}
\vspace*{-25bp}
{\color{LetterPink}\hrule height 3.5bp}\vspace{5bp}
\noindent{\sffamily\bfseries\fontsize{17bp}{21bp}\selectfont Nature Cardiovascular Research\par}
\vspace{37bp}
\noindent{\LetterType}\hfill{\LetterIdentifier}\par
\vspace{3bp}\hrule height .35bp\vspace{5bp}
\noindent{\bfseries\fontsize{26bp}{32bp}\selectfont 资料索引中的记录关联、版本标记\par\noindent 与来源定位\par}
\vspace{23bp}\hrule height .35bp\vspace{6bp}
\noindent\begin{minipage}[t]{.317\textwidth}
  \vspace{0pt}\sffamily\fontsize{8bp}{11.3bp}\selectfont
  收稿日期\enspace 待填写\par\vspace{4bp}\hrule height .3bp\vspace{4bp}
  接受日期\enspace 待填写\par\vspace{4bp}\hrule height .3bp\vspace{4bp}
  在线发表\enspace 待填写\par\vspace{4bp}\hrule height .3bp\vspace{4bp}
  非官方中文样板
\end{minipage}\hfill
\begin{minipage}[t]{.655\textwidth}
  \vspace{0pt}\setlength{\parindent}{0pt}
  {\sffamily\bfseries\fontsize{9bp}{12.1bp}\selectfont
  示例作者甲\textsuperscript{1}、示例作者乙\textsuperscript{2}、示例作者丙\textsuperscript{1,2}、
  示例作者丁\textsuperscript{1}、示例作者戊\textsuperscript{2}、示例作者己\textsuperscript{2}\par}
  \vspace{6bp}\hrule height .35bp\vspace{7bp}
  {\fontsize{11bp}{15.2bp}\selectfont 资料索引把条目名称、版本和来源定位信息汇集在一起。条目名称可以改动，索引仍需保留从每条记录返回相应原件的路径。编目时可把用于浏览的短名称与用于定位的完整题名分开，再在详页显示原件的版本说明。索引顺序调整后，读者仍可以沿用原来的定位方式。这里用原创的中性场景演示中文研究简报：图展示一条记录如何关联索引和详页，表将字段及其含义并列，附录演示报告栏目中的状态与说明。所有信息均为版式占位，正式译稿按所给原件填写。\par}
\end{minipage}
\par\vspace{13bp}

\begin{multicols}{2}
\noindent 索引记录可以由短名称、详细题名、版本标记与来源位置组成。短名称便于浏览，详细题名用于区分相似条目。若只保留短名称，两个材料可能呈现相同的索引词，却分别指向不同的原件。为避免混淆，每条记录都应在其详情中保留可核查的来源\LetterCite{1}。

新版本可以沿用原来的名称。如果索引员确知记录之间存在版本关系，可以增加指向旧版的链接，同时保留两份材料各自的版本标记。只知道新旧先后时，索引仅记录可确定的先后顺序。若缺少发布日期，索引可直接记为待核实，之后有依据时再更新。

同一条记录也可以出现在不同分类下。记录正文保留在一处，各分类入口各自使用指向详情的链接。图~1用可编辑的矢量线条表示这些关联。原始材料从索引中移除前，应核对仍指向它的分类入口，以及其中是否还有可用的版本说明。

索引还可以附一份编目说明，交代名称、版本及来源位置的填写约定。遇到未确定的字段时，应保留待核实状态；后续填充时注明信息来自何处，便于使用者根据原件复核\LetterCite{2}。

编目说明也可以列出字段次序。例如，若同一条记录有多个标题形式，索引显示读者通常检索的形式，详细名称仍在对应记录中。每种标题形式与详页的关联均可通过编辑视图查看。

纸质索引页与电子索引可以共享记录的编号。若一条记录同时呈现两个版本，带卷页的正式版本应保留这些字段，其他版本则保留各自原件所载的出版状态。索引中的字段空缺以清楚的状态展示。
\end{multicols}
\vfill\hrule height .35bp\vspace{6bp}
\noindent{\sffamily\fontsize{7.65bp}{10bp}\selectfont
  \textsuperscript{1}第一示例单位（中性占位）。\enspace
  \textsuperscript{2}第二示例单位（中性占位）。\enspace 作者贡献及通信信息按所给原稿填写。\par}

\clearpage
\begin{multicols}{2}
从索引进入详页后，读者可以按记录的原顺序阅读字段。若某字段同时包含中文名称和原文标识，应让两部分可分别识别。对具有歧义的短语可在详页中增加说明，索引列表仍使用通常的简明名称。

在纸面记录中，链接可用编号表示。编号要与详页现有标识配套，且不同组的编号应能区分。附表可作为资料的一部分，顺序位置另行记载。索引只存储必要的定位字段，附表中的完整内容留在详页。
\end{multicols}
\noindent\begin{minipage}{\textwidth}
\centering
\begin{tikzpicture}[font=\sffamily\small,>=Stealth,
  entry/.style={draw=LetterBlue,fill=LetterBlue!7,rounded corners=2bp,align=center,
    text width=88bp,minimum height=33bp,inner sep=8bp},
  detail/.style={entry,draw=LetterPink!85!black,fill=LetterPink!6}]
  \node[entry] (groupa) at (0,2.25) {按名称索引};
  \node[entry] (groupb) at (0,-2.25) {按版本索引};
  \node[entry] (record) at (5.5,0) {共同记录};
  \node[detail] (sheet) at (11,0) {来源详页};
  \draw[->,thick] (groupa.east)--(record.west);
  \draw[->,thick] (groupb.east)--(record.west);
  \draw[->,thick] (record.east)--(sheet.west);
\end{tikzpicture}\par\vspace{6bp}
\raggedright\sffamily\fontsize{7.8bp}{10.25bp}\selectfont
\textbf{图~1｜索引与来源位置的关联。}两个索引入口分别按名称和版本组织，均指向同一记录。共同记录提供来源详页的定位信息。箭头表示索引跳转，节点、颜色和连线为原创版式占位。\par
\end{minipage}\par\medskip
\begin{multicols}{2}
\section*{索引字段的维护}
字段名称应与其详细说明对应。如果整个索引采用同一种日期格式，新条目也应采用该格式。记录同时保留原来的版本标识，方便核对原件的页次和版次。对格式有特殊要求的字段，可在附录中另行说明。

修改索引顺序时，应复核记录对应的版本。目录可以按主题调整，来源定位应随记录保留。若多个条目被收入同一卷，索引还可列出卷内位置，便于在不同媒介之间查找对应条目。

\subsection*{资料范围}
有些附表仅出现在说明文字中，原件未随材料提供。索引可以保留这一关系，并说明目前可定位的材料范围。后来获得附表时，可更新访问说明并补充链接，同时保留最初的条目顺序。

\subsection*{引用与回查}
在索引正文中，数字上标可以链接本文件的占位文献。点击编号后，读者可以查看演示书目中的对应条目。正式稿替换占位书目时，需连同链接目标与引用编号一并替换。

\subsection*{记录位置的变化}
详页移动或重新编排后，应核对现有索引入口。若材料存在多个排版版本，可在详页同时标示数字页次、出版页码和相应版本，将定位从单一页数变为可复查的一组字段。

核对记录时，可先按索引顺序逐条定位，再从详页回查对应入口。对于含附表或补充材料的记录，可依次核对正文引用、附表标题与附表所在位置。记录边界变化时应更新定位字段。
\end{multicols}

\clearpage\thispagestyle{letterreport}
\section*{扩展数据与报告摘要}
\noindent{\sffamily\bfseries 扩展数据表~1｜索引字段示例}\par\vspace{4bp}
\begingroup\sffamily\fontsize{8bp}{11bp}\selectfont
\renewcommand{\arraystretch}{1.35}
\noindent\begin{tabularx}{\textwidth}{@{}p{.23\textwidth}p{.25\textwidth}X@{}}
  \toprule 字段 & 状态 & 说明 \\\midrule
  条目名称 & 示例条目 & 演示索引主名称与详页之间的对应。 \\
  版本标记 & 待填写 & 保留原件所载的版本标记。 \\
  来源位置 & 待填写 & 可填写卷内页码或材料已载明的标识。 \\
  附表信息 & 待核实 & 记录实际随稿提供的材料及缺项。 \\\bottomrule
\end{tabularx}\par
\noindent 上表列举了四种版式占位字段。正式扩展表按原稿保留列名、各行、统计条件及完整脚注。\par
\endgroup

\medskip
\ReportHeading{报告栏目与状态}
\noindent{\sffamily\fontsize{8.5bp}{11.6bp}\selectfont
  \(\square\) 来源字段待填写\qquad \(\checkmark\) 示例状态已标记\par
  选择符号、作者已填写的内容、未填写的栏目提示和不适用状态，应逐项和原件对应。\par}
\noindent\textbf{示例已填说明}\quad 参见本模板的索引字段示例，检查其中的名称、版本及来源位置三栏。\par
\noindent\textbf{待填写栏目}\quad 正式稿在核对实际所给附表后填写访问位置。\par
\ReportHeading{资料和声明}
\noindent 这一栏演示如何紧跟栏目标题列出访问条件、作者贡献或利益冲突说明。正式译稿应按原文保留各独立声明及其主语，许可信息需核对所给版本。\par
\noindent\textbf{材料范围}\quad 中性演示包括索引字段示例和链接图。实际译稿可在对应栏目列出原件提供的附件。\par
\noindent\textbf{许可字段}\quad 按原件逐字核对许可证标识、第三方素材说明及适用范围。\par

\begin{multicols}{2}
\SampleBibliography
\section*{原文信息}
\noindent 作者与英文题名\quad 按所给原文填写。\par
\noindent 卷期、页码和 DOI\quad 按所给原文填写。\par
\noindent 样板中的书目、图、表和报告状态均为原创中性占位。\par
\medskip\noindent 仅供学术交流、出版物版权归原作者及出版商所有。\par
\end{multicols}
\end{document}
